\chapter{The problem we are actually trying to solve}
\label{sec:problem}

Every project in our workspace trains something: a reinforcement-learning
(RL) agent through RLlib (Ray's RL library), a Hamiltonian
neural network on trajectory data (a network whose architecture builds in
energy conservation), a normalizing flow that preconditions a sampler (a
neural map that reshapes a hard posterior into an easier one), a
forecasting model with several competing losses. Each of these is taught
properly where it matters (Chapter~\ref{sec:workloads} devotes a section
to each); for now they are simply the customers. Each of these training procedures exposes ten to twenty
knobs, and the knobs come in four kinds. The kinds are worth meeting one
at a time, because every method in this document has to handle all four,
and more than one otherwise-good method stumbles on exactly one of them.

A \emph{continuous} knob takes any value in an interval. The learning
rate is the canonical case: $3 \times 10^{-4}$ is a legal setting, so is
$2.9 \times 10^{-4}$, and so is everything between them. Entropy
coefficients and integrator step sizes are the same kind. Continuous
knobs are the friendliest kind, because ``nearby'' means what it seems
to mean: settings close together usually behave similarly, and every
model-based method in this document leans on that.

An \emph{ordinal} knob is an integer with a natural order. Batch size
runs 128, 256, 512; network width and the number of integrator steps
inside a sampler are the same kind. The order carries real information, since 256 genuinely
sits between 128 and 512 and performance usually moves through it on
the way from one to the other, but the grid is coarse, and treating
the values as just another continuum papers over the coarseness. How
much that ordering is worth respecting returns as a concrete design
choice in Chapter~\ref{sec:bayesian}.

A \emph{categorical} knob is a choice among labels with no ``between''
at all. There is no activation function halfway from $\tanh$ to ReLU,
and no gradual path between using spectral normalization (a training
stabilizer we will meet again in the running example) and not using
it. Methods built on continuity must do something special here, and
what they do (overlap kernels, per-category bandits) occupies real
space in Chapters~\ref{sec:bayesian} and~\ref{sec:population}.

A \emph{conditional} knob exists only when another knob takes a
particular value. Learning-rate annealing is the standing instance: if
annealing is switched off, the final learning-rate multiplier is not
merely irrelevant, it is undefined, and a tuner that keeps proposing
values for it is spending budget on configurations that do not differ.
(Off-policy algorithms supply the same structure elsewhere in our
workloads: the temperature of a soft target update is undefined when
hard updates are selected.)
Conditional structure is why realistic search spaces are trees rather
than boxes, and it decides real verdicts later; it is much of the case
for one of Chapter~\ref{sec:bayesian}'s searchers.

Choosing these values well is routinely worth more than
the next clever idea, and choosing them badly can make a sound method look
broken. \citet{eimer2023hyperparameters} swept eleven hyperparameters of
PPO, the same reinforcement-learning algorithm our running example will
tune below, across standard environments and found that almost every one of them had a
large effect on final return in at least one environment, including knobs that
most papers never tune, such as the clipping range. The same study found that
the worst value of a swept hyperparameter landed within one standard deviation
of the best value in only 7 of 126 sweeps. Tuning is not optional in our line
of work; the only question is how much compute and attention it consumes.

\section{A running example, and the loop every method lives in}
\label{sec:running}

Abstract definitions slide off the mind unless they have something to stick
to, so we fix a concrete case now and return to it in every section. Our
running example is the one our own cluster runs most often: tuning PPO,
the workhorse reinforcement-learning algorithm of our RLlib stack, as it
learns a locomotion task in Brax, a suite of fast physics-simulated
environments.

Since PPO's knobs will follow us through the whole document, one
plain-language pass through the algorithm is worth its space; nothing
deeper than this pass is ever assumed. The agent holds two neural
networks. The \emph{policy} maps what the agent observes to a
distribution over actions; the \emph{value network} predicts how much
reward is still to come from a given situation. Training alternates two
activities: the agent acts in the simulated environment for a while,
collecting observations, actions, and rewards, and then the optimizer
takes gradient steps that make the actions which worked out better than
the value network expected (the \emph{advantage}) more probable, with a
clip that keeps the updated policy close to the policy that collected
the data. The name PPO, proximal policy optimization, is a description:
optimize the policy, but stay proximal. A trained agent's score, its
\emph{return}, is the total reward it accumulates over an episode of
play.

The full search space, taken from the experiments behind BG-PBT, the
population-based method Chapter~\ref{sec:population} examines in
detail, has fifteen
dimensions: nine training hyperparameters (learning rate, discount, entropy
coefficient, unroll length, reward scaling, batch size, update epochs,
the generalized-advantage-estimation parameter GAE
$\lambda$, clipping $\epsilon$) and six architecture parameters (width,
depth, and a spectral-normalization flag, separately for the policy and
value networks) \citep{wan2022bgpbt}.

\subsection{Reading the fifteen knobs}

Fifteen numbers is a lot to hold, so read them in groups, by what each
group controls. The aim is not to memorize the list but to see that
each entry is a lever a practitioner can explain, not a mystery
constant.

Three knobs govern the size of an update. The learning rate scales how
far every gradient step moves the network, and it is the most
notorious hyperparameter in deep learning for a reason: too large
diverges, too small never arrives, and the usable range spans decades.
PPO's clipping $\epsilon$ bounds how far one update may move the
\emph{policy} away from the policy that collected the current data, a
trust leash rather than a step size. Reward scaling sets the units of
the training signal itself, and every gradient inherits those units.

Two knobs decide what the future is worth. The discount $\gamma$
weights a reward $t$ steps ahead by $\gamma^{t}$, so it is exactly a
time-preference parameter, and the conventional $\gamma = 0.99$ gives
an effective horizon of about $1/(1-\gamma) = 100$ steps. GAE's
$\lambda$ (generalized advantage estimation) sets how far into the
observed future the advantage estimate
looks, trading the bias of short lookaheads against the variance of
long ones.

One knob buys experimentation. The entropy coefficient pays the policy
a bonus for staying random, so that it keeps trying actions it would
otherwise have abandoned. Too little and the agent commits early to a
mediocre habit; too much and it never commits at all.

Three knobs meter the data per update: the unroll length (how many
environment steps are collected before an update), the batch size (how
many of those steps each gradient estimate averages over), and the
number of update epochs (how many passes the optimizer makes over the
same collected data before discarding it).

The remaining six are the architecture: width and depth for the policy
network and again for the value network, capacity in four integers,
plus a spectral-normalization flag for each, a stabilizer and our
standing example of a categorical knob.

None of these levers acts alone. The size of a stable gradient step
depends on how much data estimated it, so batch size and learning rate
tend to move together; a longer unroll changes what $\gamma$ and
$\lambda$ can see; extra capacity changes how much restraining the
entropy bonus and the clip must do. This coupling is why turning one
knob at a time misleads, and why methods that model the joint surface
(Chapter~\ref{sec:bayesian}) earn their complexity. Whenever a figure
needs a space we can draw, we will slice out the two knobs everyone
recognizes: the learning rate, searched on a log scale from $10^{-5}$
to $10^{-2}$, and the entropy coefficient.

\subsection{Expensive, noisy, and without gradients}

One evaluation of one configuration means training an agent for up to
150 million environment steps, hours of wall-clock on several
accelerators. Everything else in this document is downstream of that
number: a hundred-trial study is hundreds of GPU-hours, a line item
someone has to defend, and Section~\ref{sec:costmodel} prices it out.

The resulting score is noisy. Each training run reports the mean
evaluated return over a handful of episodes, and two agents differing only in
random seed can land far apart. When the formal statement below takes
an expectation over seeds, this spread is what the expectation is
pricing.

And the score arrives without gradients in the knobs. Backpropagation
supplies derivatives of the loss with respect to the network's
weights, not with respect to the batch size that fed it or the
discount that shaped its targets; between $\config$ and the final
return stands an entire non-differentiable training procedure. Nothing
like gradient descent is available in the space we are actually
searching.

Expensive, noisy, and gradient-free: those three properties are the
whole reason this document exists.

A second, smaller case will carry the multi-objective thread: a Hamiltonian
neural network from our physics projects, where one training run is minutes
rather than hours, and where nobody can say in advance how many units of long-rollout
energy drift one unit of short-horizon prediction error is worth. We will
meet it properly in Chapters~\ref{sec:moo} and~\ref{sec:workloads}.

\subsection{The loop, and the words for its parts}

Every method in this document, from grid search to the most elaborate
population scheme, is an instance of one loop, and it is worth drawing once
so that later chapters can point at its parts
(Figure~\ref{fig:loop}). A \emph{tuner} holds two responsibilities that we
will keep separate throughout, because the literature develops them
separately and our own design (Chapter~\ref{sec:architecture}) implements
them behind separate interfaces. The \emph{searcher} decides which
configuration to try next. The \emph{scheduler} decides how long each trial
runs, which trials to stop early, and which paused trials to resume. The
executor trains the model and streams back intermediate scores; everything
observed lands in a store that the tuner (and, later, other studies) can
consult. Grid search is this loop with a trivial searcher and no scheduler;
Bayesian optimization upgrades the searcher
(Chapter~\ref{sec:bayesian}); successive halving upgrades the scheduler
(Chapter~\ref{sec:multifidelity}); population-based training entangles the
two on purpose (Chapter~\ref{sec:population}).

\begin{figure}[t]
\centering
\begin{tikzpicture}[
  box/.style={draw, rounded corners, align=center, inner sep=5pt, font=\small},
  arr/.style={-{Latex}, thick}]
\node[box, minimum width=46mm, minimum height=16mm] (tuner) at (0,0)
  {\textbf{tuner}\\[1pt] searcher: which $\config$ next?\\ scheduler: how long? stop? resume?};
\node[box, minimum width=38mm, minimum height=16mm] (exec) at (8.8,0)
  {\textbf{executor}\\[1pt] trains $A(\config; s)$\\ at fidelity $r$};
\draw[arr] ([yshift=3.5mm]tuner.east) -- node[above, font=\footnotesize] {trial $(\config, s, r)$} ([yshift=3.5mm]exec.west);
\draw[arr] ([yshift=-3.5mm]exec.west) -- node[below, font=\footnotesize] {curve, score, cost} ([yshift=-3.5mm]tuner.east);
\node[box, minimum width=118mm, minimum height=9mm] (store) at (4.4,-2.8)
  {\textbf{run store}: configurations, seeds, learning curves, costs, diagnostics};
\draw[arr] (exec.south) -- node[right, font=\footnotesize] {records} (exec.south |- store.north);
\draw[arr] (store.north -| tuner.south) -- node[left, font=\footnotesize] {informs} (tuner.south);
\node[font=\small\itshape, above=1mm of tuner.north] {proposes};
\node[font=\small\itshape, above=1mm of exec.north] {pays the GPU bill};
\end{tikzpicture}
\caption{The loop every tuning method instantiates. The searcher and the
scheduler are separate decisions: one picks configurations, the other
allocates training budget across them. Everything observed flows into a
store, which is what later lets a new study start smarter than its
predecessor (Chapter~\ref{sec:transfer}).}
\label{fig:loop}
\end{figure}

Three words from the figure recur on nearly every page, so we fix them now.
A \emph{trial} is a single training run with a fixed configuration (and, we
will insist later, a recorded seed). The \emph{incumbent} is the best
configuration found so far. The \emph{fidelity} of a partial evaluation is
its resource level: epochs trained, environment steps consumed, chain length
run. In the running example, a full-fidelity trial is 150 million steps; a
one-tenth-fidelity trial of the same configuration costs one tenth of the
GPU bill and tells us something, but not everything, about the full run.

To see the loop turn once, walk a young study through it. The searcher,
knowing nothing yet, proposes a first configuration: learning rate
$3 \times 10^{-4}$, width 256, annealing off. The scheduler grants
it a tenth of full fidelity, 15 million steps. The executor trains,
streaming back a learning curve, and the store records the
configuration, the seed, the curve, and the bill. At the next decision
point the scheduler looks at that curve and either buys the trial more
steps or stops payment, and the searcher, now holding one completed
record rather than none, proposes a second configuration a little less
blindly. Every chapter that follows upgrades one arrow of
Figure~\ref{fig:loop} and leaves the rest of the loop alone.

\section{The formal problem, and where it bends}
\label{sec:formal}

The textbook formulation is simple, and we will need it on the page because
every method chapter refines some part of it. A training algorithm $A$ maps a
configuration $\config \in \cspace$ and a random seed $s$ to a trained model,
and a cost function $c$ scores that model on held-out data, so we want
\begin{equation}
\config^{\star} \in \arg\min_{\config \in \cspace}
\; \E_{s}\!\left[ c\!\left( A(\config; s) \right) \right],
\label{eq:hpo}
\end{equation}
where the expectation runs over seeds (and, when relevant, over draws of the
training data or of environment instances). Read the pieces against the
running example: $\cspace$ is the fifteen-dimensional product of intervals
(learning rate), integer ranges (width, batch size), and finite sets (the
spectral-normalization flag), sometimes with conditional branches; $A$ is
the machine-learning training loop as a whole, here PPO and everything
around it; $c$ is the negative of the evaluated return; the expectation
over $s$ is the part beginners drop and RL practitioners regret dropping.
Each evaluation of the objective means training a model, so evaluations are
expensive, noisy, and yield no gradients with respect to $\config$. This is
black-box optimization under noise, and Chapter~\ref{sec:bayesian} describes
the machinery built for exactly this setting.

The expectation in \eqref{eq:hpo} deserves one concrete reading before
it starts doing work. Suppose, schematically, that configuration A
trains to evaluated returns of 900 and 300 on two seeds, while
configuration B trains to 640 and 660. A leaderboard that remembers
best single runs crowns A, whose 900 is the best number anyone has
seen. The expectation averages instead, 600 against 650, advantage B,
and that is the right verdict for anyone who must ship one
configuration and then live with whatever seed the production run
draws. Optimizing the average of a noisy objective and optimizing its
luckiest observed value are different problems, and \eqref{eq:hpo}
commits us, deliberately, to the first.

Formulation~\eqref{eq:hpo} bends in four places for our workloads, and each
bend picks out a branch of the literature. The four bends are worth reading
slowly: they are the map of Chapters~\ref{sec:modelfree}
through~\ref{sec:transfer}, and each one will reappear as the opening
problem of its chapter.

\subsection{Bend one: the seed is part of the problem}

The expectation over seeds is not a nuisance term; in reinforcement
learning it is often the dominant source of variation.
\citet{eimer2023hyperparameters} observed single seeds crashing while sibling
seeds of the same configuration finished well, and found that a configuration
selected on one seed can sit far below average on others. Concretely: if we
tune the running example on seed 1 and the winning configuration happened to
draw a lucky initialization, we have optimized for the luck, not the
configuration, and the ``win'' evaporates on seeds 2 through 5. Their remedy,
which we adopt in Chapter~\ref{sec:roadmap}, is to treat the seed budget as
part of the tuning protocol: evaluate candidates on several training seeds,
and report final results on test seeds that the tuner never saw. The
structure is the same train/test split that protects supervised learning
from overfitting, applied one level up, to the tuner itself. The same
discipline appears in the sampler world, where a single chain can look
excellent by luck.

\subsection{Bend two: the best configuration drifts during training}

The best configuration is not constant in time.
Learning-rate schedules exist because the loss surface a network sees at the
end of training is not the one it saw at the start, and in reinforcement
learning the data distribution itself shifts as the policy improves: the
running example's agent collects its own training data, so better play
produces different data, which changes what the next gradient step should
be. \citet{eimer2023hyperparameters} formalize the distinction as algorithm
configuration, which seeks one fixed $\config^{\star}$, versus dynamic
algorithm configuration, which seeks a policy mapping training state to
configuration; landscape studies that refit the response surface at
several points during RL training find the good region genuinely drifting
\citep{mohan2023landscapes}. Population-based training and its descendants
(Chapter~\ref{sec:population}) live in the dynamic setting: they output a
schedule of hyperparameters rather than a point, which is precisely why they
appeal for RL and why they need more parallel compute than pointwise methods.

\subsection{Bend three: several costs at once}

We often care about several costs at once. A Hamiltonian network
trades prediction error against violation of energy conservation; a trading
model trades expected return against drawdown; any model on shared hardware
trades quality against training or inference cost. When the exchange rate
between objectives is unknown in advance, the honest target is the Pareto
front, the set of configurations that cannot be improved in one objective
without paying in another; fixing weights up front would answer by
assumption the question the study exists to ask.
Chapter~\ref{sec:moo} develops this properly,
including the geometry of why the obvious fix (just add the objectives with
weights) quietly discards solutions.

\subsection{Bend four: cheap evidence is biased evidence}

A partial evaluation is cheap and informative. A learning curve at
epoch ten says a great deal about epoch one hundred; a short MCMC run says
something (though, as Chapter~\ref{sec:workloads} discusses, not everything)
about a long one. Methods that exploit partial evaluations, called
multi-fidelity methods, are the single largest lever for cutting tuning cost,
routinely saving an order of magnitude of compute over full-length
evaluations (Chapter~\ref{sec:multifidelity}). The catch, previewed here
because it shapes everything in that section: cheap evidence is biased
evidence, and the bias has a sign. Configurations that invest early (large
learning rates) look better at low fidelity than configurations that invest
late (small learning rates, heavy regularization), so a scheduler that
trusts early scores too much systematically kills the slow starters. Keeping
that bias in view is most of the art.

\section{What ``fast and cheap'' should mean}
\label{sec:costmodel}

The phrase ``low computation cost'' hides three different quantities, and a
tuner can be good at one while being poor at another. Before comparing any
two methods we need to know which bill we are trying to shrink.

The first is the compute consumed by the trials themselves: the sum of
GPU-hours over every training run the tuner launches, including the runs it
kills early. This is where essentially all the money goes, and it is the
quantity multi-fidelity methods attack. In the running example, a 100-trial
study at full fidelity is hundreds of GPU-hours; every other cost in this
subsection is rounding error next to it. A useful way to compare tuners is
the total trial compute needed to reach a target quality with, say,
probability one half, estimated over repeated tuning runs.

The second is wall-clock time to a good configuration, which depends on how
well the tuner exploits parallel workers. A method that needs its evaluations
one after another (classical Bayesian optimization in its sequential form)
can lose to a dumber parallel method simply by leaving GPUs idle;
asynchronous variants exist precisely to close that gap. The
arithmetic separating the two bills is worth one line: a
sixty-four-trial study at full fidelity burns the same GPU-hours on
one worker as on eight, but the calendar time differs by a factor of
eight, and a tuner that can only propose sequentially cannot use the
other seven workers at all.
\citet{wan2022bgpbt} report that a sequential Bayesian optimization baseline
needed roughly six times more resources than their population method to reach
comparable quality on Brax tasks, and more still in wall-clock terms, because
the population trains in parallel while sequential optimization cannot. The
distinction matters to us operationally: cluster GPU-hours and a researcher's
calendar week are different budgets, held by different people.

The third is the tuner's own overhead: fitting whatever internal model
the tuner carries, scoring candidates with it, moving checkpoints. This
is usually negligible next to
training. \citet{eimer2023hyperparameters} measured minutes of optimizer
overhead against days of trial compute in their RL study, with only their
most elaborate population method (BG-PBT) reaching about two hours of overhead at
their largest budget of 64 runs. Overhead only starts to matter when trials
themselves are short, in the regime of small Hamiltonian networks or quick
supervised probes, where a heavyweight distributed stack can dominate the
cost of the work. This is worth remembering when we weigh Ray against
lighter execution backends in Chapter~\ref{sec:architecture}.

\section{What this document does not cover}

Three neighboring topics are out of scope, with pointers for the curious.
Gradient-based hyperparameter optimization differentiates the validation loss
through the training procedure itself, either exactly on small problems or by
implicit differentiation at scale; \citet{lorraine2020millions} tune millions
of regularization hyperparameters this way. It is powerful when the
hyperparameters enter the training loss smoothly, but it requires
intrusive access to the training loop and does not apply to discrete
choices, so we treat it as a specialist tool rather than product scope.

Learned optimizers and meta-gradient methods that adapt hyperparameters inside
a single run, such as the self-tuning actor-critic of
\citet{zahavy2020selftuning}, sit inside the training algorithm rather than
outside it; they compete with outer-loop tuning and we note in
Chapter~\ref{sec:workloads} where internal adaptation should simply be
switched on and left alone.

Full neural architecture search with weight
sharing is a large field of its own; Chapter~\ref{sec:workloads} takes from
it only the pieces that have survived independent evaluation and that fit a
tuning product, and explains why we do not recommend building a
differentiable-NAS capability.
